\documentclass[../DoAn.tex]{subfiles}
\begin{document}

Chương này trình bày bài toán chấm công, lý do đồ án xử lý khuôn mặt ngay
trên vi điều khiển, mục tiêu, phạm vi và định hướng giải pháp.

\section{Đặt vấn đề}
\label{section:1.1}

Trong doanh nghiệp, dữ liệu chấm công đi thẳng vào tiền lương. Lượt vào ra
của nhân viên được quy đổi thành ngày công theo ca làm, ngày công là căn cứ
tính lương, và từ lương mới tính thuế thu nhập cá nhân và bảo hiểm xã hội.
Một lượt chấm sai hoặc bị mất làm sai phiếu lương của người đó, và bộ phận
nhân sự phải lần ngược từ phiếu lương về từng lượt chấm để sửa.

Các máy chấm công hiện nay chủ yếu dùng thẻ từ, vân tay hoặc khuôn mặt. Thẻ
từ rẻ nhưng chỉ chứng minh ai đang cầm thẻ, nên người này có thể quẹt hộ
người kia~\cite{jain2004}. Vân tay gắn với từng người nhưng buộc mọi nhân
viên chạm vào cùng một cảm biến, và vân tay của lao động chân tay thường bị
mòn hoặc xước nên máy khó đọc~\cite{jain2004}. Nhận diện khuôn mặt không cần
vật mang và không cần chạm vào máy, nhưng kéo theo hai vấn đề khác.

Vấn đề đầu tiên là giả mạo. Nếu máy chỉ so khớp khuôn mặt, một tấm ảnh in
hay một đoạn video phát trên điện thoại cũng chấm công được thay cho người
trong ảnh. Cách tấn công này gọi là tấn công trình diện, được mô tả trong bộ
tiêu chuẩn ISO/IEC 30107~\cite{iso30107}, và phát hiện khuôn mặt giả đã
thành một hướng nghiên cứu riêng của thị giác máy tính~\cite{yu2023fas}.

Vấn đề còn lại là quyền riêng tư. Dữ liệu sinh trắc học nằm trong danh mục
dữ liệu cá nhân nhạy cảm của Nghị định 356/2025/NĐ-CP~\cite{nd3562025}, và
Luật Bảo vệ dữ liệu cá nhân có hiệu lực từ ngày 01/01/2026 có một điều riêng
cho loại dữ liệu này~\cite{luat912025}. Nếu ảnh khuôn mặt được gửi về máy chủ
để nhận diện, doanh nghiệp phải bảo vệ nó cả trên đường truyền lẫn ở nơi lưu.
Nếu máy chấm công tự nhận diện, ảnh không cần rời khỏi máy.

Mô hình học sâu cho khuôn mặt thường chạy trên máy tính nhúng có hệ điều hành
đầy đủ hoặc bộ tăng tốc riêng, nên giá mỗi máy cao hơn nhiều so với một vi
điều khiển. TinyML nghiên cứu cách chạy mô hình học máy trên vi
điều khiển~\cite{warden2019,david2021}. Đồ án chọn ESP32-S3, có hai lõi Xtensa
LX7 tốc độ tối đa 240~MHz, 512~KB SRAM và tập lệnh vector dùng cho tính toán
mạng nơ-ron~\cite{esp32s3}. Khó khăn nằm ở chỗ ba mô hình phát hiện, chống giả
mạo và nhận diện khuôn mặt sau khi lượng tử hóa vẫn nặng 2.116~KB, hơn bốn lần
SRAM của chip, trong khi người chấm công chỉ đứng trước máy khoảng 1--2 giây.

Phía sau máy chấm công còn cả phần nghiệp vụ. Lượt chấm phải được tính thành
ngày công và bảng lương, nhân viên cần tra được công, phép, lương của mình và
gửi đơn trên điện thoại, người quản trị cần biết máy nào đang mất kết nối. Khi
mạng đứt, máy vẫn phải nhận diện được người và giữ lại lượt chấm để gửi sau.

\section{Mục tiêu và phạm vi đề tài}
\label{section:1.2}

Đồ án xây dựng hệ thống chấm công nhận diện khuôn mặt chạy trên ESP32-S3,
gồm cả phần quản trị nhân sự nhận dữ liệu từ máy. Sản phẩm gồm:
(i)~máy chấm công phát hiện, chống giả mạo và nhận diện khuôn mặt ngay trên
chip, cho đăng ký khuôn mặt tại máy, tự thức khi có người lại gần và vẫn chấm
công khi mất mạng; (ii)~hệ thống quản trị nhân sự trên web với hồ sơ nhân
viên, cơ cấu tổ chức, ca làm, bảng công, nghỉ phép, kỳ lương, phiếu lương và
báo cáo; (iii)~cổng nhân viên trên điện thoại, ở dạng ứng dụng web PWA và ứng
dụng Android, để xem công, xin nghỉ và nhận phiếu lương; (iv)~chức năng quản
lý máy chấm công từ xa gồm duyệt máy mới, theo dõi trạng thái và cập nhật
qua OTA.

Máy chấm công phải chấm được mà không cần tiếp xúc, chặn được ảnh in và video
phát lại, và không gửi ảnh khuôn mặt lên máy chủ. Phía máy chủ phải tự đi từ
lượt chấm tới ngày công và phiếu lương, có tính thuế thu nhập cá nhân và bảo
hiểm xã hội, để nhân viên tự tra cứu trên điện thoại thay vì hỏi bộ phận nhân
sự.

Đồ án chỉ xét tấn công bằng ảnh in và video phát lại. Mặt nạ 3D có chiều sâu
thật, nên mô hình đầu vào 80$\times$80 chạy trên vi điều khiển không đọc được
dấu vết của nó; rủi ro này được ghi nhận và không xử lý trong đồ án. Máy chấm
công được thiết kế để đặt trong nhà, nơi có người qua lại. Phần nhân sự dừng
ở chấm công, nghỉ phép và tính lương, không gồm tuyển dụng, đánh giá năng lực
và đào tạo.

\section{Định hướng giải pháp}
\label{section:1.3}

Hệ thống chia thành ba phần: máy chấm công, máy chủ và ứng dụng cho người dùng.

Máy chấm công dùng module ESP32-S3 có 16~MB flash và 8~MB PSRAM, camera
OV5640, màn hình cảm ứng 4~inch, cảm biến khoảng cách VL53L1X, loa và servo gạt
thanh chắn, gắn trên một mạch đế thiết kế bằng KiCad. Cảm biến thấy có người
thì máy thức dậy. Mỗi khuôn mặt đi qua ba mô hình nối tiếp trên chip: YuNet tìm
khuôn mặt và năm điểm mốc~\cite{yunet2023}, MiniFASNet phân biệt mặt thật với
mặt giả~\cite{minifas}, MobileFaceNet tạo vector đặc trưng 512
chiều~\cite{mobilefacenet2018} để so với các mẫu lưu trên máy. Các mô hình được
lượng tử hóa chủ yếu về số nguyên 8~bit và chạy bằng thư viện
ESP-DL~\cite{espdl}. Đo trên board, một lượt qua cả ba mô hình mất 406,2~ms khi
camera đang xem trước, nhanh hơn 9,5 lần so với cùng bộ mô hình chạy bằng
TFLite Micro.

Lượt chấm được ghi vào flash của máy trước, có mạng mới gửi lên máy chủ qua
MQTT trên TLS. Mỗi lượt mang một mã duy nhất theo từng máy nên gửi lại nhiều lần
cũng chỉ thành một bản ghi. Máy chủ dùng NestJS và PostgreSQL để tính ngày công,
bảng công và bảng lương, và gửi bản cập nhật firmware và mô hình xuống máy qua
OTA. Ảnh khuôn mặt không rời máy chấm công; vector đặc trưng trên máy chủ được
mã hóa khi lưu.

Người dùng làm việc qua bảng điều khiển web dành cho bộ phận nhân sự và quản
trị, và cổng nhân viên trên điện thoại. Các giao diện dùng chung một API và cập
nhật ngay khi dữ liệu thay đổi.

Đóng góp chính của đồ án là đưa trọn chuỗi phát hiện, chống giả mạo và nhận diện
khuôn mặt lên một vi điều khiển và đo đạc trên board thật, trong đó có so sánh
ESP-DL với TFLite Micro và đánh giá chống giả mạo trên ảnh do chính camera OV5640
chụp. Đồ án cũng xây dựng đường dữ liệu không mất lượt chấm khi mất mạng, và một
hệ quản trị nhân sự chạy trên máy chủ thật từ tháng 9/2026.

\section{Bố cục đồ án}
\label{section:1.4}

Phần còn lại của đồ án được tổ chức như sau.
Chương~\ref{chapter:Related_works} trình bày cơ sở lý thuyết: các giải pháp chấm
công, TinyML, ba bài toán phát hiện, chống giả mạo và nhận diện khuôn mặt, lượng
tử hóa mô hình, nền tảng nhúng, truyền thông IoT, công nghệ máy chủ và web, nghiệp
vụ nhân sự và các công trình liên quan.
Chương~\ref{chapter:Design} phân tích yêu cầu và trình bày thiết kế phần cứng,
pipeline AI, firmware, bản tin MQTT, cơ sở dữ liệu, API và giao diện.
Chương~\ref{chapter:Implementation} trình bày việc lắp ráp phần cứng, huấn luyện
và triển khai mô hình, xây dựng firmware, backend, các phân hệ nhân sự, giao diện
web, ứng dụng Android và triển khai hệ thống.
Chương~\ref{chapter:Experiment} đo độ trễ, bộ nhớ và độ chính xác trên board, so
sánh ESP-DL với TFLite Micro, đánh giá chống giả mạo, tỉ lệ FAR và FRR, kiểm thử
backend, kiểm thử tải và kiểm thử toàn hệ thống, rồi tổng hợp các giải pháp và
đóng góp nổi bật.
Phần Kết luận tóm tắt kết quả đạt được và đề xuất hướng phát triển.

\end{document}
